\documentclass[10pt,a4paper]{amsart}
\usepackage[margin=19mm]{geometry}
\usepackage{amsmath,amssymb,mathtools}
\usepackage[scr=rsfs,cal=euler]{mathalfa}
\usepackage{xfrac}
\usepackage[unicode,hidelinks,bookmarks=false]{hyperref}
\newcommand{\ie}{i.e.\ }
\newcommand{\cf}{cf.\ }
\newcommand{\resp}{respectively\ }
\newcommand{\Subsection}{\subsection}
\providecommand{\nobreakdash}{\nobreak\hbox{-}}
\setlength{\emergencystretch}{2em}
\multlinegap0pt
\usepackage{bm}
\usepackage{bbm}
\usepackage{dsfont}
\usepackage{xspace}
\usepackage{cancel}
\usepackage{mathtools}
\usepackage{amsthm}
\makeatletter
\@ifundefined{theorem}{\newtheorem{theorem}{Theorem}}{}
\@ifundefined{lemma}{\newtheorem{lemma}[theorem]{Lemma}}{}
\@ifundefined{proposition}{\newtheorem{proposition}[theorem]{Proposition}}{}
\makeatother
\title[The Monge--Amp\`ere equation on graphs]{The Monge--Amp\`ere equation on graphs}
\author{Ahmed Alkhozaae, Julio D. Rossi, Aelson Sobral, Jos\'e Miguel Urbano}
\date{Source: arXiv:2608.22580v1 (CC BY 4.0)}
\begin{document}
\maketitle

\noindent\emph{Source and licence.} The Monge--Amp\`ere equation on graphs. Version arXiv:2608.22580v1 (CC BY 4.0), distributed under the Creative Commons Attribution 4.0 International licence, as recorded in the arXiv licence metadata for this exact version and shown on the abstract page https://arxiv.org/abs/2608.22580v1. Full rights notice: licenses/arxiv-2608.22580-CC-BY-4.0.txt.\par\smallskip

\noindent\textit{Editorial scope.} Counted unit: the Theorem whose source label is \texttt{thm:comparison-principle} in the article (source0.tex lines 764--791, source label \texttt{thm:comparison-principle}), delivered together with the article's own complete original proof of it (source0.tex lines 793--939, one \texttt{proof} environment of 147 lines). No mathematics is written, repaired or re-derived: statement and proof are the article's own text line for line. Required context retained uncounted: lemma:restrict-set-alpha-2 (lines 545--550); prop:propagation-to-boundary (lines 686--696). Every definition, display and label of those lines is retained unchanged. Omitted from this package and not redistributed: 1--544, 551--685, 697--763, 792--792, 940--2440. Nothing inside a retained excerpt is omitted or altered: every retained body is delivered exactly as the article's lines are written, so no omission receipt of any kind accompanies this package. Citations: the retained excerpts carry no citation command at all, so no bibliography entry of the article is transcribed and no reference is redistributed. Reference adaptation: none was needed and none was made; every reference inside the delivered unit resolves inside it, so no unresolved reference or citation is produced. Printed slips kept literal: no printed word, symbol, hypothesis or number of the counted statement or of its proof was altered. Layout and class: the article's own class is \texttt{amsart} with packages including amsmath, amssymb, amsthm, mathtools, mathrsfs, bbm, graphicx, tikz, pstricks, pst-grad, pst-plot, grffile, verbatim, xcolor; the wrapper is the lane's pinned \texttt{amsart} editorial wrapper at 10pt on A4, which restates the theorem-like environments the retained text needs and adds the article's own macro definitions that the retained text needs (none), with extra wrapper packages (bm, bbm, dsfont, xspace, cancel, mathtools). The delivered numbering is the unit's own. Assets: no retained span contains a figure environment, a graphics-inclusion command, a PDF-page inclusion, a table or a TikZ picture; no image file is copied into this package, the package declares no asset dependency and no image file is redistributed with it. Licence: version arXiv:2608.22580v1 of this article is distributed under the Creative Commons Attribution 4.0 International licence, as recorded in the arXiv licence metadata for this exact version and shown on the abstract page https://arxiv.org/abs/2608.22580v1. A full rights notice is delivered at licenses/arxiv-2608.22580-CC-BY-4.0.txt. Characters: every retained excerpt is pure ASCII, so the delivered engine, XeLaTeX with its default Latin Modern text font, reports no missing character.
\begin{lemma}\label{lemma:restrict-set-alpha-2}
Let $u\colon \mathcal{X} \to \mathbb{R}$ and let $f\colon \mathcal X\setminus \mathcal O\to (0,\infty)$. For every $x\in\mathcal X\setminus\mathcal O$, there exists $\alpha^*\in\mathcal A_{n/2}$ such that
\[
    F[u](x)=\mathcal Q_u(x;\alpha^*).
\]
\end{lemma}

\begin{proposition}\label{prop:propagation-to-boundary}
Let \(u,v\colon\mathcal X\to \mathbb R\) satisfy \(u\le v\) in \(\mathcal X\) and
\[
    u(x_0)=v(x_0)
\]
for some \(x_0\in \mathcal X\). Assume moreover that
\[
    H_i[u](x_0)=H_i[v](x_0)\qquad \text{for every } i\in\{1,\dots,n/2\}.
\]
Then \(u=v\) on \(\mathcal N_{x_0}\).
\end{proposition}

\begin{theorem}\label{thm:comparison-principle}
Let \(f\colon\mathcal X\setminus\mathcal O
\to(0,\infty)\), and let \(u,v\colon\mathcal X\to\mathbb R\) satisfy
\begin{equation}\label{eq:bvp-sub}
    \begin{cases}
        \displaystyle
        u(x) \leq F[u](x),
        & x\in\mathcal X\setminus\mathcal O,\\[3ex]
        u(x)\leq g(x),
        & x\in\mathcal O,
    \end{cases}
\end{equation}
and
\begin{equation}\label{eq:bvp-super}
    \begin{cases}
        \displaystyle
        v(x) \geq F[v](x),
        & x\in\mathcal X\setminus\mathcal O,\\[3ex]
        v(x)\geq g(x),
        & x\in\mathcal O.
    \end{cases}
\end{equation}
Then
\[
    u\leq v
    \qquad\text{in }\mathcal X.
\]
\end{theorem}

\begin{proof}
Set \( m\coloneqq n/2\).  The monotonicity of the operator \(F\) follows from the monotonicity of each \(H_i\), and moreover
\begin{equation}\label{eq:F-translation}
    F[z+c](x)=F[z](x)+c.
\end{equation}
We first prove that
\begin{equation}\label{eq:F-strictly-below-H1}
    F[z](x)<H_1[z](x),
\end{equation}
for every \(z\colon\mathcal X\to\mathbb R\) and every
\(x\in\mathcal X\setminus\mathcal O\). Indeed, write
\[
    h_i\coloneqq H_i[z](x),
    \qquad i=1,\dots,m.
\]
Since \(h_1\leq\cdots\leq h_m\) and \(f(x)>0\), there exists a unique
\(t<h_1\) such that
\[
    \prod_{i=1}^{m}(h_i-t)=f(x).
\]
Define
\[
    \alpha_i
    \coloneqq
    \frac{f(x)^{1/m}}{h_i-t},
    \qquad i=1,\dots,m.
\]
Then \(\alpha_i>0\) and
\[
    \prod_{i=1}^{m}\alpha_i
    =
    \frac{f(x)}
    {\prod_{i=1}^{m}(h_i-t)}
    =1,
\]
so \(\alpha\in\mathcal A_m\). Moreover,
\[
    \alpha_i(h_i-t)=f(x)^{1/m},
\]
and therefore
\[
\mathcal{Q}_z(x;\alpha) = t+ \frac{\sum_{i=1}^{m}\alpha_i(h_i-t) -m f(x)^{1/m}}{\sum_{i=1}^{m}\alpha_i}=t.
\]
It follows that
\[
    F[z](x)\leq t<h_1=H_1[z](x),
\]
which proves \eqref{eq:F-strictly-below-H1}.

Suppose, seeking a contradiction, that
\[
    M
    \coloneqq
    \max_{x\in\mathcal X}\bigl(u(x)-v(x)\bigr)>0.
\]
Define \(A \coloneqq \left\{x\in\mathcal X \colon u(x)-v(x)=M\right\}\) and observe that since \(u\leq g\leq v\) on \(\mathcal O\), we have \(A\subseteq\mathcal X\setminus\mathcal O\). Setting \(w\coloneqq v+M\), we obtain \(u\leq w\) in \(\mathcal X\), and
\[
    u(x)=w(x)
    \qquad\text{for every }x\in A.
\]

Fix \(x\in A\). Using the subsolution inequality for \(u\), the
monotonicity and translation invariance of \(F\), and the supersolution inequality for \(v\), we obtain
\[
\begin{aligned}
    u(x) \leq F[u](x) \leq F[w](x) = F[v](x)+M \leq v(x)+M=u(x).
\end{aligned}
\]
In particular
\begin{equation}\label{eq:F-contact}
    F[u](x)=F[w](x).
\end{equation}
By Lemma \ref{lemma:restrict-set-alpha-2}, there exists
\(\alpha^*\in\mathcal A_m\) such that \(F[w](x)=\mathcal{Q}_w(x;\alpha^*)\). Using \eqref{eq:F-contact} and the definition of \(F[u]\), we get
\[
    \mathcal{Q}_w(x;\alpha^*)
    =
    F[w](x)
    =
    F[u](x)
    \leq
    \mathcal{Q}_u(x;\alpha^*).
\]
On the other hand, \(u\leq w\) implies \(\mathcal{Q}_u(x;\alpha^*)\leq \mathcal{Q}_w(x;\alpha^*)\), and hence
\[
    \mathcal{Q}_u(x;\alpha^*)=\mathcal{Q}_w(x;\alpha^*).
\]
Consequently,
\[
    0
    =
    \mathcal{Q}_w(x;\alpha^*)-\mathcal{Q}_u(x;\alpha^*)
    =
    \frac{
        \displaystyle
        \sum_{i=1}^{m}\alpha_i^*
        \bigl(H_i[w](x)-H_i[u](x)\bigr)
    }{
        \displaystyle\sum_{i=1}^{m}\alpha_i^*
    }.
\]
Every term in the numerator is nonnegative, and every
\(\alpha_i^*\) is positive. Therefore,
\[
    H_i[u](x)=H_i[w](x),
    \qquad i=1,\dots,m.
\]

Since \(u\leq w\), \(u(x)=w(x)\), and all the corresponding \(H_i\) agree at \(x\), Proposition \ref{prop:propagation-to-boundary} gives \(u=w\) on \(\mathcal N_x\), and so
\begin{equation}\label{eq:contact-set-closed}
    \mathcal N_x\subseteq A
    \qquad\text{for every }x\in A.
\end{equation}
Since \(A\) is finite and nonempty, choose
\(\widehat x\in A\) such that
\[
    u(\widehat x)=\max_{x\in A}u(x).
\]
By \eqref{eq:contact-set-closed}, every neighbor of \(\widehat x\)
belongs to \(A\). Therefore,
\[
    u(y)\leq u(\widehat x)
    \qquad\text{for every }y\in\mathcal N_{\widehat x},
\]
which implies
\[
    H_1[u](\widehat x)\leq u(\widehat x).
\]
However, the subsolution inequality and
\eqref{eq:F-strictly-below-H1} give
\[
    u(\widehat x)
    \leq
    F[u](\widehat x)
    <
    H_1[u](\widehat x)
    \leq
    u(\widehat x),
\]
which is impossible.

We conclude that
\[
    \max_{x\in\mathcal X}\bigl(u(x)-v(x)\bigr)\leq0,
\]
and therefore \(u\leq v\) in \(\mathcal X\).
\end{proof}

\end{document}
